%% ============================================================
%%  CHAPITRE 3 — Analyse et Conception du Système
%% ============================================================

\section*{Introduction}

Ce chapitre présente l'analyse et la conception détaillée du système. Nous y définissons
les fondements du DevSecOps et le principe Shift-Left, décrivons l'architecture globale en
six couches, l'architecture physique sur Google Cloud Platform, l'architecture logique
Kubernetes, les diagrammes UML de conception (séquence, activité, déploiement), le stack
technique de l'application cible, et les choix technologiques justifiés.


\section{Définition du DevSecOps et Principe Shift-Left}

\textbf{DevSecOps} est une extension de la philosophie DevOps qui intègre la sécurité
(\textbf{Sec}) comme responsabilité partagée, continue et automatisée dans le cycle de
développement. L'idée fondamentale est le \textbf{Shift-Left}~: déplacer les contrôles
de sécurité le plus tôt possible dans le cycle, là où les corrections sont les moins coûteuses.

\begin{figure}[H]
    \centering
    \fbox{\parbox{0.85\textwidth}{\centering
        \vspace{2.5cm}
        Espace réservé pour la Figure~\ref{fig:shiftleft}\\
        \textbf{Principe Shift-Left de la sécurité dans le cycle DevSecOps}\\[6pt]
        \footnotesize Plan $\to$ Code $\to$ Build $\to$ Test $\to$ Release $\to$ Deploy $\to$ Monitor\\
        Sec Review | SAST/Lint | SCA/Trivy | DAST/ZAP | cosign | K8s Sec | SIEM/Wazuh
        \vspace{2.5cm}
    }}
    \caption{Principe Shift-Left de la sécurité dans le cycle DevSecOps}
    \label{fig:shiftleft}
\end{figure}


\section{Les Huit Principes Directeurs du Projet}

Le tableau~\ref{tab:principes} présente les huit principes directeurs qui guident la conception
de la solution~:

\begin{table}[H]
\centering
\setlength{\arrayrulewidth}{0.8pt}
\setlength{\dashlinedash}{3pt}
\setlength{\dashlinegap}{2pt}
\renewcommand{\arraystretch}{1.4}
\caption{Huit principes directeurs de la plateforme DevSecOps}
\label{tab:principes}
\begin{tabularx}{\textwidth}{|>{\centering\arraybackslash}X|>{\centering\arraybackslash}X|>{\centering\arraybackslash}X|}
\hline
\tableheader \textbf{\#} & \textbf{Principe} & \textbf{Implémentation} \\
\hline
1 & Tout est Code & Terraform pour l'infra, YAML pour K8s, tout versionné dans Git \\ \hdashline
2 & Sécurité Shift-Left & Trivy + SonarCloud dès le build, avant tout déploiement \\ \hdashline
3 & Keyless First & Workload Identity Federation --- zéro clé JSON statique \\ \hdashline
4 & Zero Trust Réseau & NetworkPolicies default-deny, flux explicitement autorisés \\ \hdashline
5 & GitOps & Argo CD comme seul moteur de réconciliation du cluster \\ \hdashline
6 & Tests au plus proche du réel & Playwright E2E dans le cluster, PostSync hooks \\ \hdashline
7 & Observabilité Continue & Prometheus + Grafana + Wazuh en permanence \\ \hdashline
8 & Preuves Auditables & Chaque outil génère des rapports uploadés comme artefacts CI \\
\hline
\end{tabularx}
\end{table}


\section{Architecture Globale en Six Couches}

La solution s'articule autour de \textbf{six couches complémentaires} illustrées dans
la figure~\ref{fig:couches}~:

\begin{figure}[H]
    \centering
    \fbox{\parbox{0.85\textwidth}{\centering
        \vspace{3cm}
        Espace réservé pour la Figure~\ref{fig:couches}\\
        \textbf{Architecture en six couches de la plateforme DevSecOps}\\[6pt]
        \footnotesize
        Couche 6~: Sécurité Runtime (Wazuh, NetworkPolicies, Workload Identity)\\
        Couche 5~: Observabilité (Prometheus, Grafana, nginx-exporter)\\
        Couche 4~: Orchestration (GKE Autopilot, Namespaces, HPA, Ingress)\\
        Couche 3~: Livraison (GitHub Actions, Argo CD, Artifact Registry)\\
        Couche 2~: Sécurité Applicative (Trivy, OWASP ZAP, SonarCloud, Playwright)\\
        Couche 1~: Infrastructure (Terraform, GCP VPC, GKE, IAM, Artifact Registry)
        \vspace{3cm}
    }}
    \caption{Architecture en six couches de la plateforme DevSecOps}
    \label{fig:couches}
\end{figure}


\section{Architecture Physique GCP}

L'infrastructure est déployée sur le projet GCP \texttt{pfe-esprit-489411} dans la région
\texttt{europe-west1}. Elle comprend les composants suivants~:

\begin{itemize}
  \item \textbf{VPC \texttt{devops-vpc}}~: réseau dédié avec subnet \texttt{devops-subnet}~;
        plages secondaires pour les nœuds (\texttt{10.10.0.0/20}), pods (\texttt{10.20.0.0/16})
        et services (\texttt{10.30.0.0/20})~;
  \item \textbf{Cluster GKE Autopilot \texttt{devops-cluster}}~: mode géré, patches de
        sécurité automatiques, facturation à la ressource pod~;
  \item \textbf{Artifact Registry}~: stockage des images Docker de l'application et des tests~;
  \item \textbf{VM Compute Engine}~: hébergement du manager Wazuh SIEM
        (\texttt{e2-medium}, IP \texttt{10.10.0.2}).
\end{itemize}

\begin{figure}[H]
    \centering
    \fbox{\parbox{0.85\textwidth}{\centering
        \vspace{3cm}
        Espace réservé pour la Figure~\ref{fig:archi_physique}\\
        \textbf{Architecture physique GCP}\\[6pt]
        \footnotesize GCP Project $\to$ VPC devops-vpc $\to$ Subnet devops-subnet\\
        GKE Autopilot Cluster | Artifact Registry | Wazuh VM (e2-medium)\\
        Workload Identity Federation $\leftrightarrow$ GitHub Actions
        \vspace{3cm}
    }}
    \caption{Architecture physique de l'infrastructure GCP}
    \label{fig:archi_physique}
\end{figure}


\section{Architecture Logique Kubernetes}

Le cluster est segmenté en \textbf{cinq namespaces fonctionnels} avec des responsabilités
strictement séparées, comme présenté dans le tableau~\ref{tab:namespaces}~:

\begin{table}[H]
\centering
\setlength{\arrayrulewidth}{0.8pt}
\setlength{\dashlinedash}{3pt}
\setlength{\dashlinegap}{2pt}
\renewcommand{\arraystretch}{1.4}
\caption{Organisation en namespaces Kubernetes}
\label{tab:namespaces}
\begin{tabularx}{\textwidth}{|>{\centering\arraybackslash}X|>{\centering\arraybackslash}X|>{\centering\arraybackslash}X|}
\hline
\tableheader \textbf{Namespace} & \textbf{Rôle} & \textbf{Ressources principales} \\
\hline
\texttt{app} & Workload applicatif & Deployment, Service, Ingress, HPA, PostSync Job \\ \hdashline
\texttt{testing} & Validation E2E & Playwright Job, ServiceAccount \\ \hdashline
\texttt{observability} & Collecte de métriques & Prometheus, Grafana, ServiceMonitor \\ \hdashline
\texttt{security} & Politiques et SIEM & NetworkPolicies, Wazuh Agent DaemonSet \\ \hdashline
\texttt{argocd} & Moteur GitOps & Argo CD Applications \\
\hline
\end{tabularx}
\end{table}

\begin{figure}[H]
    \centering
    \fbox{\parbox{0.85\textwidth}{\centering
        \vspace{3cm}
        Espace réservé pour la Figure~\ref{fig:archi_k8s}\\
        \textbf{Diagramme d'architecture logique Kubernetes}\\[6pt]
        \footnotesize 5 namespaces~: app, testing, observability, security, argocd\\
        Flux réseau~: testing $\to$ app (TCP 80), observability $\to$ app (TCP 9113),\\
        security $\to$ Wazuh VM (TCP 1514), Ingress $\to$ app (TCP 80)
        \vspace{3cm}
    }}
    \caption{Architecture logique du cluster Kubernetes}
    \label{fig:archi_k8s}
\end{figure}


\section{Diagrammes de Conception}

\subsection{Diagramme de Séquence~: Pipeline CI/CD Complet}

Le diagramme de séquence suivant modélise le flux complet du pipeline CI/CD, depuis le
push du développeur jusqu'au déploiement GitOps et les tests E2E.

\begin{figure}[H]
    \centering
    \fbox{\parbox{0.85\textwidth}{\centering
        \vspace{4cm}
        Espace réservé pour la Figure~\ref{fig:seq_pipeline}\\
        \textbf{Diagramme de séquence~: Pipeline CI/CD complet}\\[6pt]
        \footnotesize Développeur $\to$ GitHub~: push commit\\
        GitHub $\to$ Actions~: déclencher pipeline\\
        Actions $\to$ SonarCloud~: scan SAST\\
        Actions $\to$ Trivy~: scan FS (SCA)\\
        Actions $\to$ Artifact Registry~: push image Docker\\
        Actions $\to$ Trivy~: scan image\\
        Actions $\to$ Git~: update K8s tags\\
        Argo CD $\to$ Git~: detect change\\
        Argo CD $\to$ GKE~: sync ressources\\
        Argo CD $\to$ Playwright~: PostSync E2E tests\\
        Actions $\to$ ZAP~: scan DAST
        \vspace{4cm}
    }}
    \caption{Diagramme de séquence du pipeline CI/CD complet}
    \label{fig:seq_pipeline}
\end{figure}

\subsection{Diagramme de Séquence~: Authentification Workload Identity Federation}

\begin{figure}[H]
    \centering
    \fbox{\parbox{0.85\textwidth}{\centering
        \vspace{3cm}
        Espace réservé pour la Figure~\ref{fig:seq_wif}\\
        \textbf{Diagramme de séquence~: Authentification WIF}\\[6pt]
        \footnotesize GitHub Actions $\to$ GitHub OIDC~: request JWT token\\
        GitHub Actions $\to$ GCP WIF Pool~: present JWT\\
        GCP WIF $\to$ GCP IAM~: verify issuer, repo, branch\\
        GCP IAM $\to$ GitHub Actions~: temporary credentials (1h)\\
        GitHub Actions $\to$ GCP APIs~: authenticated calls
        \vspace{3cm}
    }}
    \caption{Diagramme de séquence de l'authentification Workload Identity Federation}
    \label{fig:seq_wif}
\end{figure}

\subsection{Diagramme d'Activité~: Flux GitOps avec Argo CD}

\begin{figure}[H]
    \centering
    \fbox{\parbox{0.85\textwidth}{\centering
        \vspace{3cm}
        Espace réservé pour la Figure~\ref{fig:act_gitops}\\
        \textbf{Diagramme d'activité~: Flux GitOps Argo CD}\\[6pt]
        \footnotesize [Commit dans Git] $\to$ Argo CD détecte le changement $\to$\\
        Comparer état désiré vs état réel $\to$ [Si différence] Synchroniser $\to$\\
        Appliquer les manifests K8s $\to$ Exécuter PostSync hooks (Playwright) $\to$\\
        [Si drift détecté] Self-heal automatique $\to$ État synchronisé
        \vspace{3cm}
    }}
    \caption{Diagramme d'activité du flux GitOps avec Argo CD}
    \label{fig:act_gitops}
\end{figure}

\subsection{Diagramme de Déploiement}

\begin{figure}[H]
    \centering
    \fbox{\parbox{0.85\textwidth}{\centering
        \vspace{3.5cm}
        Espace réservé pour la Figure~\ref{fig:deployment_diagram}\\
        \textbf{Diagramme de déploiement UML}\\[6pt]
        \footnotesize
        \textbf{Nœud GitHub}~: Actions Runner, OIDC Provider\\
        \textbf{Nœud GCP}~: VPC devops-vpc\\
        \quad \textbf{GKE Cluster}~: ns app (Deployment + nginx-exporter sidecar),\\
        \quad\quad ns testing (Playwright Job), ns observability (Prometheus, Grafana),\\
        \quad\quad ns security (Wazuh Agent), ns argocd (Argo CD)\\
        \quad \textbf{VM Wazuh}~: Wazuh Manager (TCP 1514)\\
        \quad \textbf{Artifact Registry}~: Docker images\\
        \textbf{Nœud Externe}~: Navigateur utilisateur $\to$ Ingress GCE
        \vspace{3.5cm}
    }}
    \caption{Diagramme de déploiement de la plateforme}
    \label{fig:deployment_diagram}
\end{figure}


\section{L'Application Cible~: Todo App React}

\subsection{Description Fonctionnelle}

L'application déployée est une \textbf{Todo App} avec un portail de connexion, développée
avec React 18 et TypeScript. Cette application a été conçue comme \textbf{support de
démonstration} pour valider l'ensemble de la chaîne DevSecOps. En effet, l'équipe QA
utilise au quotidien le framework BDD interne \textbf{T4T} pour tester les applications RH
de production~; la Todo App permet de concevoir et de valider les pipelines DevSecOps dans un
environnement maîtrisé, avant de transposer ces pratiques aux projets réels de l'équipe.

Bien que fonctionnellement simple, l'application sert de support idéal car elle possède~:

\begin{itemize}
  \item Une \textbf{authentification} (cible pour les tests de sécurité)~;
  \item Une \textbf{interaction utilisateur} (cible pour les tests E2E et XSS)~;
  \item Une \textbf{persistance de données} via localStorage (cible pour les vulnérabilités
        de stockage)~;
  \item Une \textbf{chaîne de dépendances npm} (cible pour le SCA).
\end{itemize}

\begin{figure}[H]
    \centering
    \fbox{\parbox{0.85\textwidth}{\centering
        \vspace{2.5cm}
        Espace réservé pour la Figure~\ref{fig:app_screenshot_login}\\
        \textbf{Capture d'écran~: Page de login de la Todo App}
        \vspace{2.5cm}
    }}
    \caption{Page de login de la Todo App}
    \label{fig:app_screenshot_login}
\end{figure}

\begin{figure}[H]
    \centering
    \fbox{\parbox{0.85\textwidth}{\centering
        \vspace{2.5cm}
        Espace réservé pour la Figure~\ref{fig:app_screenshot_dashboard}\\
        \textbf{Capture d'écran~: Dashboard Todo avec liste de tâches}
        \vspace{2.5cm}
    }}
    \caption{Dashboard de gestion des tâches Todo}
    \label{fig:app_screenshot_dashboard}
\end{figure}

\subsection{Diagramme de Classes de l'Application}

\begin{figure}[H]
    \centering
    \fbox{\parbox{0.85\textwidth}{\centering
        \vspace{3cm}
        Espace réservé pour la Figure~\ref{fig:class_diagram}\\
        \textbf{Diagramme de classes de la Todo App}\\[6pt]
        \footnotesize Classes~: App, AuthContext, AuthProvider, LoginPage,\\
        TodoDashboard, TodoItem, Todo (interface)\\
        Relations~: App $\to$ AuthProvider $\to$ LoginPage / TodoDashboard\\
        TodoDashboard $\to$ TodoItem[] , TodoItem $\to$ Todo
        \vspace{3cm}
    }}
    \caption{Diagramme de classes de l'application Todo App}
    \label{fig:class_diagram}
\end{figure}

\subsection{Stack Technique}

Le tableau~\ref{tab:stack} détaille les technologies utilisées~:

\begin{table}[H]
\centering
\setlength{\arrayrulewidth}{0.8pt}
\setlength{\dashlinedash}{3pt}
\setlength{\dashlinegap}{2pt}
\renewcommand{\arraystretch}{1.4}
\caption{Stack technique de l'application}
\label{tab:stack}
\begin{tabularx}{\textwidth}{|>{\centering\arraybackslash}X|>{\centering\arraybackslash}X|>{\centering\arraybackslash}X|}
\hline
\tableheader \textbf{Couche} & \textbf{Technologie} & \textbf{Version} \\
\hline
Framework UI & React & 18.3.1 \\ \hdashline
Langage & TypeScript & 5.8.3 \\ \hdashline
Build tool & Vite & 8.0.0 \\ \hdashline
Styling & Tailwind CSS + Rosé Pine theme & 3.4.17 \\ \hdashline
Composants UI & shadcn/ui (Radix UI) & --- \\ \hdashline
Animations & Framer Motion & 12.x \\ \hdashline
Routing & React Router DOM & 6.30.1 \\ \hdashline
Gestionnaire de paquets & Bun & Latest \\ \hdashline
Serveur de production & Nginx & 1.21 (vulnérable intentionnel) \\
\hline
\end{tabularx}
\end{table}

\subsection{Vulnérabilités Intentionnelles}

\begin{warningbox}
Le composant \texttt{TodoItem} utilise intentionnellement \texttt{dangerouslySetInnerHTML}
pour rendre le texte des tâches, créant une vulnérabilité \textbf{XSS stockée} détectable
par OWASP ZAP et SonarCloud. De même, l'authentification est hardcodée côté client dans
\texttt{AuthContext.tsx}. Ces choix sont délibérés pour la démonstration des outils DevSecOps.
\end{warningbox}


\section{Choix Technologiques et Justifications}

\begin{table}[H]
\centering
\setlength{\arrayrulewidth}{0.8pt}
\setlength{\dashlinedash}{3pt}
\setlength{\dashlinegap}{2pt}
\renewcommand{\arraystretch}{1.4}
\caption{Choix technologiques et justifications}
\label{tab:choix_techno}
\begin{tabularx}{\textwidth}{|>{\centering\arraybackslash}X|>{\centering\arraybackslash}X|>{\centering\arraybackslash}X|}
\hline
\tableheader \textbf{Domaine} & \textbf{Outil choisi} & \textbf{Justification} \\
\hline
IaC & Terraform & Standard industriel, multi-cloud, état déclaratif \\ \hdashline
Orchestration & GKE Autopilot & Gestion nœuds par Google, sécurité par défaut, scaling auto \\ \hdashline
CI/CD & GitHub Actions & Natif GitHub, OIDC intégré, marketplace riche \\ \hdashline
GitOps & Argo CD & Self-heal, drift detection, UI dashboard, CNCF graduated \\ \hdashline
SAST & SonarCloud & Analyse TypeScript/React, Quality Gates, gratuit open-source \\ \hdashline
SCA + Container & Trivy & 3 modes (fs, image, config), SARIF output, open-source \\ \hdashline
DAST & OWASP ZAP & Standard industriel, baseline scan intégrable en CI \\ \hdashline
SIEM & Wazuh & Open-source, HIDS + FIM + compliance, compatible K8s \\ \hdashline
Observabilité & Prometheus + Grafana & Standard CNCF, écosystème exporters, dashboards \\ \hdashline
Tests E2E & Playwright & Multi-navigateur, conteneurisable, tags pour parallélisme \\
\hline
\end{tabularx}
\end{table}


\section{Comparaison Globale Avant/Après}

Le tableau~\ref{tab:avant_apres} offre une vue synthétique de la transformation conçue~:

\begin{table}[H]
\centering
\setlength{\arrayrulewidth}{0.8pt}
\setlength{\dashlinedash}{3pt}
\setlength{\dashlinegap}{2pt}
\renewcommand{\arraystretch}{1.4}
\caption{Comparaison des approches avant et après transformation}
\label{tab:avant_apres}
\begin{tabularx}{\textwidth}{|>{\centering\arraybackslash}X|>{\centering\arraybackslash}X|>{\centering\arraybackslash}X|}
\hline
\tableheader \textbf{Dimension} & \textbf{Avant (Jenkins)} & \textbf{Après (DevSecOps)} \\
\hline
Sécurité & Post-déploiement ou absente & Intégrée à chaque étape \\ \hdashline
Infrastructure & ClickOps manuel & Terraform IaC versionné \\ \hdashline
Déploiement & Scripts impératifs & GitOps déclaratif (Argo CD) \\ \hdashline
Tests & Unitaires seulement & Unit + E2E in-cluster + scans sécurité \\ \hdashline
Authentification & Clés JSON statiques & OIDC Workload Identity Federation \\ \hdashline
Réseau & Ouvert par défaut & Zero Trust (default-deny) \\ \hdashline
Monitoring & CPU/RAM basique & Métriques + SIEM + alertes \\ \hdashline
Audit trail & Inexistant & Chaque commit, build, déploiement tracé \\ \hdashline
Rollback & Manuel & Automatique via Git revert + Argo CD \\
\hline
\end{tabularx}
\end{table}


\section*{Conclusion}

Ce chapitre a présenté la conception détaillée du système~: une architecture DevSecOps en
six couches articulée autour de huit principes directeurs, une infrastructure GCP bien définie,
un cluster Kubernetes segmenté en cinq namespaces, et des diagrammes UML modélisant les flux
de livraison, de sécurité et de déploiement. Le chapitre suivant se consacre à l'implémentation
concrète de cette conception, organisée en quatre sprints Agile.
